\chapter{Conclusions}\label{ch:conclusions}

The three chapters at the core of this dissertation were undertaken with a shared aim: to move our understanding of attention--memory interactions from correlational observation toward a state-resolved, causally tested, and methodologically reproducible account. Together, they tell a story in which the moment-to-moment dynamics of arousal and attention --- ordinarily treated as background noise to be averaged away --- are foregrounded as legitimate targets of measurement, theory, and intervention. This concluding chapter steps back from the specifics of each contribution, synthesizes their joint implications, acknowledges their limitations, and traces out what I see as the most promising directions for future work.

\section{Synthesis of Contributions}

Chapter~\ref{ch:attending-to-remember} laid the conceptual groundwork. By reviewing recent advances in the methods and theory of attention--memory interactions, it argued that whether and how we remember is not solely a function of what was encoded but also of the preparatory states that immediately precede acts of encoding and retrieval. The chapter highlighted four developments that, in combination, are reshaping the field: a more refined neural-network account of top-down and bottom-up attention and their interaction with cognitive control; emerging evidence for rhythmic dynamics that link attention and memory at theta and alpha timescales; a sharpened ability to measure age-related changes in memory through assays of representational precision; and a growing toolkit of closed-loop paradigms that enable causal inference about attention's role in memory. Taken together, these developments motivated the empirical and methodological program that the subsequent chapters pursued.

Chapter~\ref{ch:biofeedback-framework} sketched the design space from which the empirical work was drawn. Developed early in my doctoral training as a conceptual blueprint, it laid out a taxonomy of four closed-loop pupillometry implementations differing in temporal grain (trial-by-trial vs.\ millisecond-by-millisecond) and inference strategy (fixed thresholding vs.\ individualized predictive modeling). Although the chapter never matured into a published paper, it is included here because it foreshadows --- and in places explicitly anticipates --- the methodological choices that defined the eventual PRISM paradigm.

Chapter~\ref{ch:eyeris} addressed an unglamorous but consequential prerequisite for that program: the absence of a standardized, reproducible preprocessing pipeline for pupillometry data. By introducing \texttt{eyeris} --- an R package designed around transparent, modular, and FAIR-aligned principles --- the chapter offered both a practical tool for individual researchers and a step toward bringing pupillometry preprocessing into parity with the conventions that have matured in fMRI and M/EEG. The package contributes to the broader community in three ways. First, it operationalizes a recommended preprocessing recipe in code, making the consequences of each step inspectable and adjustable rather than buried in custom scripts. Second, it generates interactive diagnostic reports that surface signal quality at the trial level, lowering the activation energy for the kind of careful data scrubbing that reliable inference requires. Third, by organizing outputs into a BIDS-like derivative structure, it positions pupillometry data for the same kinds of large-scale, multi-site analyses that have become routine in adjacent modalities. The methodological standards that \texttt{eyeris} encodes are not theoretically neutral: they reflect a stance that analytic flexibility, when unconstrained, can quietly undermine the inferences a study is meant to support, and that a transparent default workflow is one means of mitigating that risk.

Chapter~\ref{ch:prism} put both the theoretical framing and the methodological infrastructure to work. The PRISM paradigm used real-time readouts of tonic pupil diameter to detect pupil-indexed suboptimal arousal--attentional (A-A) states immediately before each retrieval attempt and, on a subset of those trials, delivered a brief visual reorienting cue intended to perturb the ongoing A-A state. Three results stand out. First, pupil-defined suboptimal states impaired both item recognition and source recollection relative to optimal-state trials, replicating and extending the readiness-to-remember framework with a within-subject, within-session causal contrast. Second, the reorienting trigger \textit{causally} rescued item memory back to optimal-state performance levels --- but not source memory --- suggesting that brief, non-semantic phasic interventions can boost the familiarity signal that underlies item recognition without thereby recovering the goal-mediated recollection of contextual detail. Third, a post hoc reclassification analysis revealed that the trigger's efficacy is bounded by both the depth and the directionality of suboptimality. Marginally suboptimal trials within a zone of recoverability benefited from the cue; deeply suboptimal trials did not. And the benefit was largest for hypoarousal (constriction-driven) departures from optimality, consistent with the ascending limb of an inverted-U arousal--performance curve, while hyperarousal (dilation-driven) departures did not benefit and may even be harmed by an A-A-boosting cue.

\section{Cross-Cutting Themes}

Three themes emerge from these chapters that I think will continue to shape the field.

\subsection*{Methods and Theory Advance Together.}

A recurring lesson is that theoretical progress in attention--memory research has been paced, in important ways, by the maturation of measurement and analysis methods. The shift from between-condition averages toward state-resolved, within-trial analyses --- so central to the readiness-to-remember framework --- was not possible until pupillometric, electrophysiological, and behavioral assays of moment-to-moment state became routine. Likewise, causal claims about preparatory states required the emergence of closed-loop paradigms that could measure those states online and act on them in the moment. The \texttt{eyeris} contribution is, in this sense, a theoretical contribution as much as a methodological one: by lowering the cost and increasing the rigor of pupillometry preprocessing, it widens the set of theoretical questions that can be addressed reliably.

\subsection*{State Matters, but the Geometry of State Matters More.}

A second theme is that ``suboptimal state'' is not a single condition but a family of states arranged along a continuum and, importantly, with a sign. The PRISM reclassification analysis made this explicit by distinguishing marginally suboptimal from deeply suboptimal departures from optimality and by separating low-arousal (constriction) from high-arousal (dilation) departures. The fact that a single, fixed-magnitude intervention produced asymmetric effects across this geometry --- rescuing some states and not others --- is itself a constraint on theory. Any future account of preparatory state and memory will need to specify not only that state matters but how its position on the inverted-U interacts with the form and magnitude of intervention. This is a richer and more empirically tractable theoretical target than ``arousal'' or ``attention'' considered as scalar quantities.

\subsection*{Pupil as Instrument, Not as Endpoint.}

Throughout this work, the pupil has served as a window onto the broader LC--NE-mediated arousal system rather than as an object of study in its own right. The accumulating evidence that pupil dynamics index a coordinated brain-wide neuromodulatory state \parencite{joshi2016relationships,joshi2020pupil} means that the pupil's value lies in its temporal resolution, non-invasiveness, and accessibility --- traits that make it a particularly tractable handle on a process that is otherwise difficult to measure continuously in awake humans. Treating the pupil as instrument rather than endpoint also clarifies how it should be deployed alongside other measures. The closed-loop paradigm developed here is, in principle, agnostic to the precise neuromodulatory mechanism the pupil reflects; what matters is that it provides a reliable, low-latency readout of preparatory state from which interventions can be triggered in real time.

\section{Limitations}

Several limitations of the work as a whole warrant acknowledgement.

First, the closed-loop study tested a single, simple intervention --- a brief flickering visual cue --- chosen for experimental tractability. Although its non-semantic character is a feature for isolating phasic arousal/attention effects from goal-content effects, it is also a constraint: it cannot speak to whether more elaborate, semantically targeted, or multimodal cues might recover source recollection in ways the present cue did not.

Second, the empirical work focused on young, healthy adults under laboratory conditions. Whether the patterns observed here generalize to older adults, to populations with attentional dysregulation (e.g., ADHD, post-concussion syndrome, early-stage neurodegenerative disease), or to ecologically valid contexts (e.g., classrooms, clinics, the workplace) is an empirical question that the present design cannot answer.

Third, pupil diameter is an indirect, non-specific proxy for LC--NE activity. While converging evidence supports its utility as a coarse continuous index, the mapping is not one-to-one, and the inferences drawn here about ``A-A state'' should be read as inferences about a pupil-indexed neurocognitive state rather than a pure neuromodulatory one.

Fourth, several of the post hoc analyses in the PRISM chapter --- in particular the directional asymmetry and depth-of-suboptimality contrasts --- were directionally consistent with theoretical predictions but did not achieve the statistical precision required for definitive adjudication. Future work, with larger samples or adaptive designs that oversample specific cells of the state continuum, will be necessary to formally test whether the trigger operates differentially across these subspaces.

Finally, \texttt{eyeris} currently supports only SR Research EyeLink \texttt{.asc} files as input. Although the underlying preprocessing principles generalize across eye-tracking hardware, extending the package to other trackers --- and adapting its default parameters to their distinct artifact profiles --- is an ongoing community effort.

\section{Future Directions}

These limitations point toward what I see as the most exciting avenues for follow-on work.

\paragraph{Tailored, adaptive interventions.} The PRISM data hint that the right intervention depends on the state being intervened upon. An A-A-boosting cue helps low-arousal suboptimal states but not high-arousal ones. A natural next step is to design adaptive closed-loop systems that read both the sign and the magnitude of departure from optimality and deploy intervention strategies accordingly --- for instance, alerting cues for deep hypoarousal states and refocusing or calming cues for hyperarousal states. Building and validating such systems will require both larger datasets and theoretically motivated, mechanistically distinct intervention libraries.

\paragraph{Recovering source memory.} The selective rescue of item but not source memory raises a focused empirical question: what intervention, delivered in the same closed-loop window, would rescue source recollection? Candidates include longer-duration or semantically richer cues, multimodal stimulation, or cues that explicitly reinstate the encoding goal context. Source-memory rescue, if achievable, would also have direct implications for dual-process theories of recognition.

\paragraph{Translational extensions.} Several populations would seem natural beneficiaries of closed-loop attention/arousal monitoring at the moment of remembering. Older adults experiencing age-related fluctuations in attentional control, individuals with ADHD or related dysregulation, and patients with mild cognitive impairment or early Alzheimer's pathology --- where LC--NE compromise is now recognized as an early feature \parencite{sara2009locus} --- are all candidate populations in which targeted, real-time intervention might mitigate state-driven memory failures. Whether such interventions can be ported out of the laboratory and into ecologically meaningful contexts, and whether they remain effective when delivered repeatedly over longer durations, are open questions.

\paragraph{Multimodal integration.} Pupillometry is a powerful but not sufficient readout of preparatory state. Combining pupil-based real-time triggering with complementary readouts --- scalp EEG markers of attention (e.g., posterior alpha power), behavioral indices of sustained attention (e.g., response-time variability), or, in carefully designed paradigms, intracranial or functional imaging measures --- would yield richer characterizations of the state space and potentially more sensitive intervention triggers. The \texttt{eyeris} infrastructure is designed to be extensible in precisely this direction.

\paragraph{Theoretical refinement of the readiness-to-remember framework.} The findings reported here also bear on the theory itself. The zone-of-recoverability account that emerged from the PRISM reclassification analysis is a fairly specific empirical claim about how the inverted-U relates to intervention efficacy. Formalizing this account --- ideally in a computational model that links pupil-indexed state, intervention magnitude, and memory outcome --- would sharpen its predictions and make them more falsifiable. Doing so will likely require integrating insights from selection-history accounts of attention, prediction-error frameworks for mnemonic salience, and dual-process models of recognition memory.

\section{Closing Remarks}

This dissertation began with a question about variability: why do we remember some moments well, and others not at all, even when the to-be-remembered content is held constant? The answer it advances is neither that variability is irreducible noise nor that it is fully under deliberate control. It is, instead, the structured product of moment-to-moment fluctuations in arousal and attention --- fluctuations that we can now measure with increasing fidelity, model with growing precision, and, as the closed-loop work here demonstrates, perturb in real time with consequences for what we are able to remember.

If there is a single takeaway from the chapters that precede this one, it is that the preparatory moments before remembering are not a passive prelude to retrieval but an active, intervenable substrate of mnemonic success. Understanding that substrate --- and learning to support it when it falters --- is, I believe, one of the most productive frontiers in contemporary memory research. The work presented here is offered as one set of steps along that frontier, with the hope that the tools, paradigms, and findings will be useful to others working toward the same goal.
